\documentclass[11pt]{article}

\usepackage{epsfig}
\usepackage{amsfonts}
\usepackage{amssymb}
\usepackage{amstext}
\usepackage{amsmath}
\usepackage{xspace}
\usepackage{theorem}
\usepackage{xcolor}
\usepackage{graphicx}
\usepackage{url}


% This is the stuff for normal spacing
\makeatletter
 \setlength{\textwidth}{6.5in}
 \setlength{\oddsidemargin}{0in}
 \setlength{\evensidemargin}{0in}
 \setlength{\topmargin}{0.25in}
 \setlength{\textheight}{8.25in}
 \setlength{\headheight}{0pt}
 \setlength{\headsep}{0pt}
 \setlength{\marginparwidth}{59pt}

 \setlength{\parindent}{0pt}
 \setlength{\parskip}{5pt plus 1pt}
 \setlength{\theorempreskipamount}{5pt plus 1pt}
 \setlength{\theorempostskipamount}{0pt}
 \setlength{\abovedisplayskip}{8pt plus 3pt minus 6pt}
\makeatother


%%%%%%%%%%%%%%%%%%%%%%%%%%%%%%%%%%%%%%%%%%%%%%%%%%%%%%%%%%%%%%%%%%%%%%%%%%%
% Document begins here %%%%%%%%%%%%%%%%%%%%%%%%%%%%%%%%%%%%%%%%%%%%%%%%%%%%
%%%%%%%%%%%%%%%%%%%%%%%%%%%%%%%%%%%%%%%%%%%%%%%%%%%%%%%%%%%%%%%%%%%%%%%%%%%

\newcommand{\headings}[3]{{\large
{\bf CSE548: Analysis of Algorithms, Fall 2026} \hfill {{\bf Release Date:} #1}\\
\rule[0.01in]{\textwidth}{0.025in}}
\begin{center} {{\bf\huge{Homework \##2}}\\{\bf( Due Date: #3 )}} \end{center}
}

\begin{document}

\headings{Oct 2}{1}{Oct 16}
%\headings{{\underline{\hspace{1.5cm}\textcolor{white}{.}}}}{{\underline{\hspace{1cm}}}}{{\underline{\hspace{1.6cm}}}}
\newcommand{\lecnum}{5}  % Lecture Number


\vspace{2in}

{\Large
\begin{center}

{\sc{Group Number:} {\underline{\hspace{1in}}}}~\\~\\

\scalebox{0.83}{
\begin{tabular}{||l|c|c||}\hline\hline
 \multicolumn{3}{||c||}{Group Members}\\ \hline
Name\hspace{3in}                              &      \hspace{0.5in}SBU ID\hspace{0.5in}     &     \% Contribution \\ \hline\hline
                                  &                 &                     \\
                                  &                 &                     \\ \hline
                                  &                 &                     \\
                                  &                 &                     \\ \hline
                                  &                 &                     \\
                                  &                 &                     \\ \hline\hline
\end{tabular}
}
\end{center}
}


\newpage

{\Large
\begin{center}

{\sc{\underline{Collaborating Groups}}}~\\
\vspace{0.25cm}
\scalebox{0.83}{
\begin{tabular}{||c|c||}\hline\hline
Group         &  Specifics \\
Number        &  (e.g., specific group member? specific task/subtask?) \\ \hline\hline
              &                                      \\ 
              &                                      \\ 
              &                                      \\ \hline
              &                                      \\ 
              &                                      \\ 
              &                                      \\ \hline
              &                                      \\ 
              &                                      \\ 
              &                                      \\ \hline
              &                                      \\ 
              &                                      \\ 
              &                                      \\ \hline
              &                                      \\ 
              &                                      \\ 
              &                                      \\ \hline\hline
\end{tabular}
}
\end{center}
}



{\Large
\begin{center}

{\sc{\underline{External Resources Used}}}~\\
\vspace{0.25cm}
\scalebox{0.83}{
\begin{tabular}{||c|c||}\hline\hline
              &  Specifics \\
              &   (e.g., cite papers, webpages, etc. and mention why each was used) \\ \hline\hline
              &                                      \\ 
    $1.$      & Labib et al., CPM 2019, Theorem 19 (Task 1 bucketing). \\ 
              &                                      \\ \hline
              &                                      \\ 
    $2.$      & OpenAI Codex: derivation, writing, and verification. \\ 
              &                                      \\ \hline
              &                                      \\ 
    $3.$      &                                      \\ 
              &                                      \\ \hline
              &                                      \\ 
    $4.$      &                                      \\ 
              &                                      \\ \hline
              &                                      \\ 
    $5.$      &                                      \\ 
              &                                      \\ \hline\hline
\end{tabular}
}
\end{center}
}



\newpage
\section*{Task 1. Minimizing the distance between two pegged frames}

Let $A_u,B_v$ be the heights at $u,v\in\{0,\ldots,n-1\}^2$, and put $N=n^2$.
We minimize the permitted separation without interpenetration; touching is
allowed.\footnote{We interpret $\alpha^p+\alpha^q$ as restricting feasible
separations. Requiring exact first contact as well can make the problem
unsolvable: for $n=1$, heights $1,6$, and $\alpha=2$, contact occurs at $7$,
which is not a sum of two nonnegative powers of $2$.}

\textbf{Reduce to integer shifts.}
Ignore zero-area boundary contacts and measure distances in cell widths.
In each coordinate, the overlapping cell offsets for displacement $t$ are
$\lfloor t\rfloor,\lceil t\rceil$, restricted to valid offsets.
Since $|t|<n$, rounding toward zero selects a valid offset. Doing this in
both coordinates keeps a nonempty subset of overlapping pairs, so cannot
increase the required separation.
Thus it suffices to use
$s\in\{-(n-1),\ldots,n-1\}^2$, with
\[
 H(s)=\max_{v-u=s}(A_u+B_v),\qquad h=\min_s H(s).
\]
\textbf{A fast feasibility test (dominance matching \cite{dominance}).}
For a threshold $d$, compute the number of violating pairs at every shift:
\[
 C_d(s)=\sum_{v-u=s}{\bf1}[A_u+B_v>d].
\]
Then $d\ge h$ if and only if some $C_d(s)=0$.
Sort the $N$ top cells by height and partition them into consecutive blocks
of at most $b$ cells, retaining their original coordinates. Handle violations
as follows.
\begin{enumerate}
\item For each block $I$, let $a_I$ be its minimum height. Form two binary
grids $X_I(u)={\bf1}[u\in I]$ and
$Y_I(v)={\bf1}[a_I+B_v>d]$. Add their cross-correlation
\[
 (X_I\star Y_I)(s)=\sum_{v-u=s}X_I(u)Y_I(v)
\]
to $C_d$, initially zero, using FFT with $X_I$ reversed in both coordinates
and both dimensions padded to at least $2n-1$. Shift $s$ uses convolution
entry $(n-1+s_1,n-1+s_2)$; extra padding entries are not legal shifts.
\item For each bottom cell $v$, binary-search the sorted top heights for the
first height greater than $d-B_v$. If it lies strictly inside a block,
explicitly add the violating pairs in that block to their shifts $v-u$.
There are at most $b$ such pairs per bottom cell.
\end{enumerate}
For a fixed $v$, the violating top heights form a suffix: whole blocks are
counted in step 1, and its only possible partial block in step 2. Hence every
violation is counted exactly once, including when heights repeat.
The test costs
\[
 O\!\left(\frac{N^2\log N}{b}+Nb+N\log N\right)
 =O(N^{3/2}\sqrt{\log N})
\]
with $b=\min\{N,\lceil\sqrt{N\log N}\rceil\}$ for $N\ge2$.
Process the convolutions one at a time to use $O(N)$ space.

\newpage
\textbf{Find the exact threshold without listing $N^2$ sums.}
Let $a_1\le\cdots\le a_N$ and $b_1\le\cdots\le b_N$ be the sorted heights.
Since $h$ is one of the sums $a_i+b_j$, search these $N$ sorted rows implicitly.
Keep an interval of remaining indices in each row and an incumbent
$z=+\infty$. Repeat:
\begin{enumerate}
\item Take each nonempty row's middle value, weighted by the number of its
remaining entries, and let $d$ be a weighted median of these values.
\item If $d$ is feasible, set $z=\min(z,d)$ and discard all entries $\ge d$.
Otherwise discard all entries $\le d$. Locate each row's new endpoint by
binary search.
\end{enumerate}
When no entries remain, return $h=z$. The invariant is that the optimum is
either the incumbent or a remaining entry. In either branch, rows whose
midpoints lie on the discarded side have at least half the total weight;
at least half of each such row is discarded. Thus at least one quarter of
all remaining entries disappear per iteration. There are $O(\log N)$ tests,
with $O(N\log N)$ additional work per iteration. Consequently,
\[
 \text{time to find }h=O(N^{3/2}\log^{3/2}N)
 =O(n^3\log^{3/2}n),\qquad \text{space}=O(n^2).
\]

\textbf{Use the restriction $\alpha^p+\alpha^q$.}
Since all heights are at least $1$, $h\ge2$. Write $p\ge q\ge0$ and define
\[
 L=\left\lceil\log_\alpha(h/2)\right\rceil,
 \qquad U=\left\lceil\log_\alpha(h-1)\right\rceil.
\]
For each integer $p\in[L,U]$, set
\[
 q_p=\begin{cases}
 0,&h-\alpha^p\le1,\\
 \left\lceil\log_\alpha(h-\alpha^p)\right\rceil,&h-\alpha^p>1.
 \end{cases}
 \qquad
 \boxed{D=\min_{L\le p\le U}\bigl(\alpha^p+\alpha^{q_p}\bigr).}
\]
This range is complete: $p<L$ gives $\alpha^p+\alpha^q<h$, whereas
$p>U$ is worse than the feasible candidate $\alpha^U+1$.
For every remaining $p$, the displayed $q_p$ is the smallest feasible
exponent and satisfies $q_p\le p$. Therefore $D$ is exactly the smallest
allowed feasible separation, not an approximation obtained by rounding pegs.
The number of candidates is at most $2+\lceil\log_\alpha2\rceil$.

For fixed $\alpha>1$, the total running time is
\[
 \boxed{O(n^3\log^{3/2}n)=o(n^4).}
\]
If $\alpha$ varies with the input, retain the additive
$O(1+\log_\alpha2)$ term; the subquartic claim is for fixed $\alpha$, not
uniform over input-dependent $\alpha$ approaching $1$.
These are arithmetic-operation bounds, with
power/logarithm evaluation counted at unit cost; exact power comparisons
resolve logarithm endpoints: $\lceil\log_\alpha x\rceil$ means the smallest
integer $k$ with $\alpha^k\ge x$, not an unchecked floating-point ceiling.
Binary convolutions compute integer counts exactly (using sufficient FFT
precision to round each coefficient with error less than $1/2$).
If a placement is desired, rerun the feasibility test at $h$ and return any
zero-count shift, which is also feasible at $D$. For $n=1$, simply take
$h=A_{00}+B_{00}$ and apply the same rounding step.

\newpage
\section*{Task 2. Multiplying $3\times3$ matrices}

\textbf{(a)} No counterexample exists for commuting scalar entries.
For any row $i$, write $(x,y,z)=(a_{i1},a_{i2},a_{i3})$ and subtract the
shared products first:
\begin{align*}
 p_{i1}-q_1&=xy+yb_{21}+xb_{12},\\
 p_{i2}-q_2&=xz+zb_{31}+xb_{13},\\
 p_{i3}-q_3&=yz+zb_{32}+yb_{23}.
\end{align*}
Substituting these into the three recovery formulas cancels every unwanted
term and gives
\[
 (c_{i1},c_{i2},c_{i3})
 =(xb_{11}+yb_{21}+zb_{31},\;
   xb_{12}+yb_{22}+zb_{32},\;
   xb_{13}+yb_{23}+zb_{33}).
\]
Thus the $3+3\cdot6=21$ scalar multiplications do compute $AB$.

\textbf{(b)} The formulas cannot be used recursively on general matrix
blocks, because their cancellations require commutativity.
For example, set all blocks of $B$ to zero, set
$A_{11}=X$, $A_{12}=Y$, and set all other blocks of $A$ to zero.
The proposed formula returns
\[
 C_{11}=-XY+YX,
\]
although the correct product is zero. Taking
\[
 X=\begin{pmatrix}0&1\\0&0\end{pmatrix},\qquad
 Y=\begin{pmatrix}0&0\\1&0\end{pmatrix}
\]
gives $YX-XY=\operatorname{diag}(-1,1)\ne0$.
Embedding these matrices into the upper-left corner of $3\times3$ blocks
gives the same failure for $n=9=3^2$.
Although the formal recurrence $T(n)=21T(n/3)+\Theta(n^2)$ would give
$\Theta(n^{\log_3 21})=o(n^{\log_2 7})$, it describes an incorrect block
algorithm. Unlike Strassen's identities, these identities do not hold over
noncommutative rings.

\section*{Task 3. Strassen under memory constraints}

\textbf{(a)} For the usual implementation that materializes temporary
matrices, the minimum asymptotic extra space is $\boxed{\Theta(n^2)}$,
excluding the input and output matrices. A constant number of
$(n/2)\times(n/2)$ buffers suffices: compute the seven products sequentially
and immediately accumulate each into its output blocks. Reuse the buffers,
so the peak auxiliary space satisfies
\[
 S(n)=S(n/2)+\Theta(n^2)=\Theta(n^2).
\]
Conversely, even one materialized half-size temporary requires
$\Omega(n^2)$ space. This bound concerns the temporary-matrix model in the
question, rather than algorithms that overwrite inputs or recompute entries.

\newpage
\textbf{(b)} Initialize $C=0$, disjoint from both inputs. Each call tries to
reserve three half-size buffers without blocking; failure releases any
partial allocation. On failure, use eight ordinary block products into
views of $C$. Each child retries allocation, exploiting later memory.

\begin{verbatim}
Accumulate(A, B, C):                 // C <- C + A B
    n <- number of rows of A
    if n = 1:
        C[0,0] <- C[0,0] + A[0,0] * B[0,0]; return
    partition A, B, C into four block views each
    buffers <- TryAllocate(3 * (n/2)^2 entries)
    if buffers succeeded:
        (X, Y, P) <- buffers
        for each of the seven rules in the table below:
            fill X and Y with the indicated operands
            set P to zero
            Accumulate(X, Y, P)
            apply the indicated signed updates to C
        free buffers
    else:
        for i, j, k in {1,2}^3:
            Accumulate(A[i,k], B[k,j], C[i,j])
\end{verbatim}

\[
\begin{array}{c|l|l|l}
 &X&Y&\text{updates after }P\leftarrow XY\\ \hline
1&A_{11}+A_{22}&B_{11}+B_{22}&C_{11}\mathrel{+}=P,\ C_{22}\mathrel{+}=P\\
2&A_{21}+A_{22}&B_{11}&C_{21}\mathrel{+}=P,\ C_{22}\mathrel{-}=P\\
3&A_{11}&B_{12}-B_{22}&C_{12}\mathrel{+}=P,\ C_{22}\mathrel{+}=P\\
4&A_{22}&B_{21}-B_{11}&C_{11}\mathrel{+}=P,\ C_{21}\mathrel{+}=P\\
5&A_{11}+A_{12}&B_{22}&C_{11}\mathrel{-}=P,\ C_{12}\mathrel{+}=P\\
6&A_{21}-A_{11}&B_{11}+B_{12}&C_{22}\mathrel{+}=P\\
7&A_{12}-A_{22}&B_{21}+B_{22}&C_{11}\mathrel{+}=P
\end{array}
\]

Both branches preserve $C\leftarrow C+AB$ by Strassen's identities or
$(AB)_{ij}=\sum_k A_{ik}B_{kj}$; induction proves correctness for any
allocation outcomes. Calls never wait and halve the size, so they terminate.
Allocated buffers remain valid until freed; input/output and stack storage
are assumed available. Block views make no copies.

If all allocations succeed, time is $\Theta(n^{\log_2 7})$; if all fail, it
is $\Theta(n^3)$ with only $O(\log n)$ stack space beyond the inputs and
output. Any mixture takes $O(n^3)$ time and at most $O(n^2)$ temporary space,
assuming each allocation attempt takes $O(n^2)$ time or less.
For other sizes, use a logical size $M=2^{\lceil\log_2 n\rceil}<2n$ and
\emph{virtual} zero padding: input views return zero outside the original
arrays, and output views ignore writes outside the original $C$. Buffer
views, when allocated, retain all their entries. Thus padding itself never
requires an extra matrix allocation, and the same bounds hold with $M=O(n)$.

\newpage
\section*{Task 4. A new transform}

Write $u=(u_0,u_1)$, with both halves of length $m=n/2$, and let $F_n$
denote the given transform. The code restores its input before returning:
the high half is restored recursively, and the low half is explicitly saved
and restored. Its output is therefore
\[
 F_n(u)=\bigl(F_m(u_0+u_1),\ F_m(u_0),\ F_m(u_1)\bigr),
 \qquad F_1(u)=u.
\]
The three blocks appear in precisely this order.

\textbf{(a)} Each output block has length
$r=3^{\log_2m}$, the value computed by the loop. Thus
\[
 n_v=3^{\log_2 n}=n^{\log_2 3},\qquad
 T(n)=3T(n/2)+\Theta(n)=\boxed{\Theta(n^{\log_2 3})}.
\]
The $O(\log n)$ loop overhead is absorbed by the linear copying work.

\textbf{(b)} The last two blocks already encode the two original halves,
so an inverse need not read the redundant first block. To make the linear
running time explicit, use offsets into the input and a single output array;
no recursive call copies a transformed block. The wrapper also supports a
product mode $d=2$, defined in part (c).

\begin{verbatim}
BackwardTransform(v, n, d = 1):
    ell <- 3^(log_2 n)              // compute once using integer powers
    if d = 2: return Product(v, 0, n, ell)
    allocate u[0 .. n-1]
    Recover(v, 0, u, 0, n, ell)
    return u

Recover(v, iv, u, iu, n, ell):
    if n = 1:
        u[iu] <- v[iv]; return
    m <- n/2; r <- ell/3
    Recover(v, iv+r,   u, iu,   m, r)
    Recover(v, iv+2*r, u, iu+m, m, r)
\end{verbatim}

For ordinary inversion, induction shows that the two calls recover $u_0$
and $u_1$ directly into the two halves of $u$, giving
\[
 B_1(n)=2B_1(n/2)+O(1)=\boxed{\Theta(n)}.
\]
This is optimal for writing $n$ coefficients. The transformed vector is
assumed valid; checking all its redundant entries is unnecessary. Literal
array-copying concatenation would instead take $\Theta(n\log n)$, and
copying transformed input blocks would cost $\Theta(n^{\log_2 3})$.
The output uses $\Theta(n)$ space and the recursion stack $O(\log n)$.

\textbf{(c)} Pad the coefficient arrays to a common power-of-two length $n$
(or $O(n)$ if padding is needed). Compute
\[
 v_A=\operatorname{ForwardTransform}(A),\quad
 v_B=\operatorname{ForwardTransform}(B),\quad
 w_i=(v_A)_i(v_B)_i.
\]
Return $\operatorname{BackwardTransform}(w,n,2)$, of length $2n-1$.

\newpage
The product decoder uses all three blocks. It returns an array of exactly
$2n-1$ coefficients; all input blocks are accessed by offsets.
\begin{verbatim}
Product(w, iw, n, ell):
    if n = 1: return [w[iw]]
    m <- n/2; r <- ell/3
    S <- Product(w, iw,     m, r)
    L <- Product(w, iw+r,   m, r)
    H <- Product(w, iw+2*r, m, r)
    allocate R[0 .. 2*n-2], initially zero
    for t <- 0 to 2*m-2:
        R[t]     <- R[t]     + L[t]
        R[m+t]   <- R[m+t]   + S[t] - L[t] - H[t]
        R[2*m+t] <- R[2*m+t] + H[t]
    return R
\end{verbatim}
The updates are additions: the shifted coefficient ranges can overlap.
For example, $A=B=1+x$ gives $w=(4,1,1)$ and product coefficients
$(1,2,1)$; the ordinary inverse alone would return $(1,1)$.

To see why the product mode works, write
$A=A_0+x^mA_1$ and $B=B_0+x^mB_1$. The three blocks of $w$ recursively
encode the products
\[
 S=(A_0+A_1)(B_0+B_1),\qquad L=A_0B_0,\qquad H=A_1B_1.
\]
Consequently the backward step reconstructs
\[
 L+x^m(S-L-H)+x^{2m}H
 =A_0B_0+x^m(A_0B_1+A_1B_0)+x^{2m}A_1B_1=AB.
\]
At $n=1$ this is just scalar multiplication, completing the induction.
Ordinary inversion alone would discard the cross terms; the product mode
uses the sum block to recover them. Equivalently, the forward transform
samples each split at $0,1,\infty$, and products require quadratic rather
than linear interpolation.

Each backward merge adds and shifts $O(n)$ coefficients, so
\[
 B_2(n)=3B_2(n/2)+\Theta(n)=\Theta(n^{\log_2 3}).
\]
The two forward transforms and the $n^{\log_2 3}$ pointwise products have
the same bound. The total time is therefore
$\boxed{\Theta(n^{\log_2 3})=o(n^2)}$.

\begin{thebibliography}{1}
\bibitem{dominance}
K. Labib, P. Uzna\'{n}ski, and D. Wolleb-Graf.
\emph{Hamming Distance Completeness}. CPM 2019, pp.\ 14:1--14:17.
Theorem 19: bucketing and convolution for less-than pattern matching.
\url{https://doi.org/10.4230/LIPIcs.CPM.2019.14}.
\end{thebibliography}

\end{document}

